\documentclass[10pt,a4paper]{amsart}
\usepackage[margin=19mm]{geometry}
\usepackage{amsmath,amssymb,mathtools}
\usepackage[scr=rsfs,cal=euler]{mathalfa}
\usepackage{xfrac}
\usepackage[unicode,hidelinks,bookmarks=false]{hyperref}
\newcommand{\ie}{i.e.\ }
\newcommand{\cf}{cf.\ }
\newcommand{\resp}{respectively\ }
\newcommand{\Subsection}{\subsection}
\providecommand{\nobreakdash}{\nobreak\hbox{-}}
\setlength{\emergencystretch}{2em}
\multlinegap0pt
\usepackage{bm}
\usepackage{bbm}
\usepackage{dsfont}
\usepackage{xspace}
\usepackage{cancel}
\usepackage{mathtools}
\usepackage{cleveref}
\usepackage[shortlabels]{enumitem}
\usepackage{amsthm}
\makeatletter
\@ifundefined{theorem}{\newtheorem{theorem}{Theorem}}{}
\@ifundefined{lemma}{\newtheorem{lemma}[theorem]{Lemma}}{}
\makeatother
\DeclareMathOperator*{\rk}{rank}
\DeclareMathOperator*{\tr}{tr}
\newcommand{\innpr}[2]{{\left\langle {#1},{#2}\right\rangle}} %inner product
\newcommand{\m}[1]{\mathbf{#1}} % General matrix notation
\renewcommand{\c}[1]{\mathcal{#1}} % Generic space notation
\title[Gaussian Optimal Transport Beyond Brenier's Theorem]{Gaussian Optimal Transport Beyond Brenier's Theorem}
\author{Ho Yun, Yoav Zemel}
\date{Source: arXiv:2512.21464v1 (CC BY 4.0)}
\begin{document}
\maketitle

\noindent\emph{Source and licence.} Gaussian Optimal Transport Beyond Brenier's Theorem. Version arXiv:2512.21464v1 (CC BY 4.0), distributed under the Creative Commons Attribution 4.0 International licence, as recorded in the arXiv licence metadata for this exact version and shown on the abstract page https://arxiv.org/abs/2512.21464v1. Full rights notice: licenses/arxiv-2512.21464-CC-BY-4.0.txt.\par\smallskip

\noindent\textit{Editorial scope.} Counted unit: the Theorem whose source label is \texttt{thm:sol:sym:inf} in the article (source0.tex lines 2284--2307, source label \texttt{thm:sol:sym:inf}), delivered together with the article's own complete original proof of it (source0.tex lines 2309--2393, one \texttt{proof} environment of 85 lines). No mathematics is written, repaired or re-derived: statement and proof are the article's own text line for line. Required context retained uncounted: lem:Schur:op (lines 418--458); eq:cond2:opt (lines 2029--2032); eq:cond:ricatti (lines 2029--2032). Every definition, display and label of those lines is retained unchanged. Omitted from this package and not redistributed: 1--417, 459--2028, 2033--2283, 2308--2308, 2394--4476. Nothing inside a retained excerpt is omitted or altered: every retained body is delivered exactly as the article's lines are written, so no omission receipt of any kind accompanies this package. Citations: the retained excerpts carry no citation command at all, so no bibliography entry of the article is transcribed and no reference is redistributed. Reference adaptation: none was needed and none was made; every reference inside the delivered unit resolves inside it, so no unresolved reference or citation is produced. Printed slips kept literal: no printed word, symbol, hypothesis or number of the counted statement or of its proof was altered. Layout and class: the article's own class is \texttt{imsart} with packages including geometry, fancyhdr, amsthm, amsmath, amsfonts, amssymb, hyperref, cleveref, mathrsfs, graphicx, float, rotating, subcaption, pdfpages; the wrapper is the lane's pinned \texttt{amsart} editorial wrapper at 10pt on A4, which restates the theorem-like environments the retained text needs and adds the article's own macro definitions that the retained text needs (\texttt{\textbackslash c}, \texttt{\textbackslash innpr}, \texttt{\textbackslash m}, \texttt{\textbackslash rk}, \texttt{\textbackslash tr}), with extra wrapper packages (bm, bbm, dsfont, xspace, cancel, mathtools, cleveref, enumitem). The delivered numbering is the unit's own. Assets: no retained span contains a figure environment, a graphics-inclusion command, a PDF-page inclusion, a table or a TikZ picture; no image file is copied into this package, the package declares no asset dependency and no image file is redistributed with it. Licence: version arXiv:2512.21464v1 of this article is distributed under the Creative Commons Attribution 4.0 International licence, as recorded in the arXiv licence metadata for this exact version and shown on the abstract page https://arxiv.org/abs/2512.21464v1. A full rights notice is delivered at licenses/arxiv-2512.21464-CC-BY-4.0.txt. Characters: every retained excerpt is pure ASCII, so the delivered engine, XeLaTeX with its default Latin Modern text font, reports no missing character.
\begin{lemma}\label{lem:Schur:op}
Let $\c{H}$ be a Hilbert space, and $\m{A}, \m{B} \in \c{B}_{0}^{+}(\c{H})$. Let $\m{G}_{11} \in \c{G}(\m{A}_{11})$ be a Green's operator, and denote its trivial extension to $\c{H}$ by
\begin{align*}
    \m{G} =
    \begin{pmatrix}
    \m{G}_{11} & \m{0} \\
    \m{0} & \m{0}
    \end{pmatrix}
    \    :
    \begin{array}{c}
    \c{H}_{1} \\
    \oplus \\
    \c{H}_{2}
    \end{array}
    \rightarrow
    \begin{array}{c}
    \c{H}_{1} \\
    \oplus \\
    \c{H}_{2}
    \end{array}.
\end{align*}
Then the following statements hold:
\begin{enumerate}[leftmargin = *]
    \item $\m{G}_{11}^{*} \in \c{B}_{0}(\c{H}_{1})$ is injective, and its extension $\m{G} \in \c{B}_{0}(\c{H})$ is a Green's operator for $\m{A}$.
    \item On $\c{H}_{1}$, $\m{G}_{11}^{*} \m{B}_{11} \m{G}_{11} = \m{G}^{*} \m{B} \m{G}$, hence $\c{N}(\m{B}_{11}^{1/2} \m{G}_{11}) = \c{N}(\m{B}^{1/2} \m{G}) \cap \c{H}_{1}$.
    \item The Schur complement $\m{B}/\m{A}$ can also be expressed as:
    \begin{equation*}
        \m{B}/\m{A} = \m{B}_{22} - [(\m{G}_{11}^{*} \m{B}_{11} \m{G}_{11})^{\dagger/2} (\m{G}_{11}^{*} \m{B}_{12})]^{*} [(\m{G}_{11}^{*} \m{B}_{11} \m{G}_{11})^{\dagger/2} (\m{G}_{11}^{*} \m{B}_{12})],
    \end{equation*}
    for any $\m{G}_{11} \in \c{G}(\m{A}_{11})$. 
    \item The operator $\m{B}^{1/2} \m{\Pi}_{\c{N}(\m{G}^{*} \m{B}^{1/2})} \m{B}^{1/2}$ is block-diagonal, where its restriction to $\c{H}_{2}$ is the $\m{A}$-Schur complement:
    \begin{equation*}
        \m{B}^{1/2} \m{\Pi}_{\c{N}(\m{G}^{*} \m{B}^{1/2})} \m{B}^{1/2} =
        \begin{cases}
            \m{0}, &\quad \text{on } \c{H}_{1}, \\
            \m{B}/\m{A}, &\quad \text{on } \c{H}_{2},
        \end{cases}
    \end{equation*}
    hence $\c{R} [(\m{B}/\m{A})^{1/2}] = \c{R}(\m{B}^{1/2}) \cap \c{H}_{2}$.
\end{enumerate}
\end{lemma}

\begin{align}
    &(\m{T} \m{A}^{1/2}) (\m{T} \m{A}^{1/2})^{*} = \m{B}, \label{eq:cond:ricatti} \\ 
    &\tr (\m{A}^{1/2} \m{T} \m{A}^{1/2}) = \tr (\m{T} \m{A}) = \tr [(\m{A}^{1/2} \m{B} \m{A}^{1/2})^{1/2}]. \label{eq:cond2:opt}
\end{align}

\begin{theorem}\label{thm:sol:sym:inf}
Let $\m{A}, \m{B} \in \c{B}_{1}^{+}(\c{H})$ be covariance operators.
Then $\m{A} \leadsto \m{B}$ if and only if the dimensional inequality holds:
\begin{align*}
    \dim [\c{R}(\m{B}^{1/2}) \cap \c{H}_{2}] \le \dim [\c{N}(\m{B}^{1/2} \m{A}^{1/2}) \cap \c{H}_{1}].
\end{align*}
If it does, any optimal pre-pushforward has the form
\begin{align}\label{eq:sol:set}
    \m{T} = 
    \begin{pmatrix}
    \m{A}_{11}^{\dagger/2} (\m{A}_{11}^{1/2} \m{B}_{11} \m{A}_{11}^{1/2})^{1/2} \m{A}_{11}^{\dagger/2} \\
    [(\m{A}_{11}^{1/2} \m{B}_{11} \m{A}_{11}^{1/2})^{\dagger/2} \m{A}_{11}^{1/2} \m{B}_{12} + \m{U}_{12} (\m{B}/\m{A})^{1/2}]^{*} \m{A}_{11}^{\dagger/2}
    \end{pmatrix}
    \    :
    \c{R}(\m{A}^{1/2}) \subset \c{H}_{1}
    \rightarrow
    \begin{array}{c}
    \c{H}_{1} \\
    \oplus \\
    \c{H}_{2}
    \end{array},
\end{align}
where $\m{U}_{12}$ is a partial isometry with initial space $\overline{\c{R}(\m{B}/\m{A})}$, and final space contained in $\c{N}(\m{B}_{11}^{1/2} \m{A}_{11}^{1/2})$.
\end{theorem}

\begin{proof}[Proof of \cref{thm:sol:sym:inf}]
\quad
\begin{itemize}[leftmargin = *]
\item [($\Rightarrow$)] Choose an optimal pre-pushforward $\m{T}: \c{R}(\m{A}^{1/2}) \subset \c{H}_{1} \rightarrow \c{H}$. Note that \eqref{eq:cond:ricatti} amounts to the fact that $\m{P} := (\m{T} \m{A}^{1/2})^{\dagger} \m{B}^{1/2}$ is a bounded operator satisfying $\m{B}^{1/2} = (\m{T} \m{A}^{1/2}) \m{P}$. Also, from \eqref{eq:cond2:opt}, we have $\m{A}^{1/2} \m{T} \m{A}^{1/2} = (\m{A}^{1/2} \m{B} \m{A}^{1/2})^{1/2}$, which is a symmetric bounded operator.

To prove the claim, it suffices to show that
\begin{equation*}
    \m{P} \m{B}^{\dagger/2} \mid_{\c{H}_{2}} : \c{R}(\m{B}^{1/2}) \cap \c{H}_{2} \rightarrow \c{N}(\m{B}^{1/2} \m{A}^{1/2}) \cap \c{H}_{1}.
\end{equation*}
is a well-defined injective operator. For well-definedness, we claim that $\m{P} \m{B}^{\dagger/2} y \in \c{N}(\m{B}^{1/2} \m{A}^{1/2}) \cap \c{H}_{1}$ for any $y \in \c{R}(\m{B}^{1/2}) \cap \c{H}_{2}$. By letting $z = \m{B}^{\dagger/2} y$ so that
\begin{align}\label{eq:injec:contra}
    y = \m{B}^{1/2} z = (\m{T} \m{A}^{1/2}) \m{P} z \in \c{H}_{2},
\end{align}
it is equivalent to showing that $\m{P} z = \m{P} \m{B}^{\dagger/2} y \in \c{N}(\m{B}^{1/2} \m{A}^{1/2}) \cap \c{H}_{1}$. For any $x \in \c{H}_{1}$,
\begin{align*}
    0 = \innpr{\m{A}^{1/2} x}{y} = \innpr{\m{A}^{1/2} x}{(\m{T} \m{A}^{1/2}) \m{P} z}
    = \innpr{x}{(\m{A}^{1/2} \m{T} \m{A}^{1/2}) \m{P} z},
\end{align*}
which establishes $(\m{A}^{1/2} \m{T} \m{A}^{1/2}) \m{P} z  = 0$, thus 
\begin{equation*}
    \m{P} z \in \c{N} (\m{A}^{1/2} \m{T} \m{A}^{1/2}) = \c{N}[(\m{A}^{1/2} \m{B} \m{A}^{1/2})^{1/2}] = \c{N}(\m{B}^{1/2} \m{A}^{1/2}).
\end{equation*}
On the other hand,
\begin{align*}
    \m{P} z \in \c{R}[(\m{T} \m{A}^{1/2})^{\dagger}] = \c{N}(\m{T} \m{A}^{1/2})^{\perp} \subset \c{N}(\m{A}^{1/2})^{\perp} = \c{H}_{1}.
\end{align*}

Finally, the injectivity of $\m{P} \m{B}^{\dagger/2}: y \mapsto \m{P} z$ is immediate from \eqref{eq:injec:contra}: $\m{P} z = 0$ if and only if $y = 0$.

\item [($\Leftarrow$)] 
Let $\m{T}: \c{R}(\m{A}^{1/2}) \subset \c{H}_{1} \rightarrow \c{H}$, and define $\m{M}_{i1} := \m{T}_{i1} \m{A}_{11}^{1/2} \in \c{B}_{\infty}(\c{H}_{1}, \c{H}_{i})$ for $i=1, 2$. Then \eqref{eq:cond:ricatti} becomes: 
\begin{equation}\label{eq:cond:ricatti:pf}
    \m{M}_{i1} \m{M}_{j1}^{*} = \m{B}_{ij}, \quad i, j = 1, 2,
\end{equation}
and \eqref{eq:cond2:opt} becomes:
\begin{equation*}
    \tr[ \m{A}_{11} \m{M}_{11}] = \tr [(\m{A}_{11}^{1/2} \m{B}_{11} \m{A}^{1/2})^{1/2}].
\end{equation*}
These two conditions determine $\m{T}_{11}$ and $\m{M}_{11}$ uniquely:
\begin{align*}
    &\m{T}_{11} = \m{S}_{\m{A}_{11} \leadsto \m{B}_{11}} = \m{A}_{11}^{\dagger/2} (\m{A}_{11}^{1/2} \m{B}_{11} \m{A}_{11}^{1/2})^{1/2} \m{A}_{11}^{\dagger/2}: \c{R}(\m{A}_{11}^{1/2}) \subset \c{H}_{1} \rightarrow \c{H}_{1}, \\
    &\m{M}_{11} = \m{T}_{11} \m{A}_{11}^{1/2} = \m{A}_{11}^{\dagger/2} (\m{A}_{11}^{1/2} \m{B}_{11} \m{A}_{11}^{1/2})^{1/2} : \c{H}_{1} \rightarrow \c{H}_{1}.
\end{align*}

Hence, it suffices to show that there exists $\m{M}_{21}$ satisfying \eqref{eq:cond:ricatti:pf} for $(i, j) =(1, 2), (2, 2)$ (as the case $(2, 1)$ is redundant):
\begin{align}
    \m{A}_{11}^{\dagger/2} (\m{A}_{11}^{1/2} \m{B}_{11} \m{A}_{11}^{1/2})^{1/2} \m{M}_{21}^{*} &= \m{B}_{12}: \c{H}_{2} \rightarrow \c{H}_{1}, \label{eq:cond:ricatti:pf1} \\
    \m{M}_{21} \m{M}_{21}^{*} &= \m{B}_{22}: \c{H}_{2} \rightarrow \c{H}_{2}. \label{eq:cond:ricatti:pf2}
\end{align}

Consider a s.p.d. block operator defined by:
\begin{align}\label{eq:comp:block}
    \m{C} := \begin{pmatrix}
        \m{C}_{11} & \m{C}_{12} \\
        \m{C}_{21} & \m{C}_{22}
    \end{pmatrix} = \begin{pmatrix}
        \m{A}_{11}^{1/2} \m{B}_{11} \m{A}_{11}^{1/2} & \m{A}_{11}^{1/2} \m{B}_{12} \\
        \m{B}_{21} \m{A}_{11}^{1/2} & \m{B}_{22}
    \end{pmatrix}
    \    :
    \begin{array}{c}
    \c{H}_{1} \\
    \oplus \\
    \c{H}_{2}
    \end{array}
    \rightarrow
    \begin{array}{c}
    \c{H}_{1} \\
    \oplus \\
    \c{H}_{2}
    \end{array}.
\end{align}
Since $\c{H}_{1} = \c{N}(\m{A}_{11}^{1/2})^{\perp}$, \eqref{eq:cond:ricatti:pf1} is equivalent to $\m{C}_{11}^{1/2} \m{M}_{21}^{*} = \m{C}_{12}$, which always have a solution characterized by
\begin{equation*}
    \m{M}_{21}^{*} = \m{C}_{11}^{\dagger/2} \m{C}_{12} \oplus \m{N}_{12}, \quad \m{N}_{12} \in \c{B}(\c{H}_{2}, \c{H}_{1}), \quad \c{R}(\m{N}_{12}) \subset \c{N}(\m{C}_{11}).
\end{equation*}
Applying this form into \eqref{eq:cond:ricatti:pf2}, we obtain from \cref{lem:Schur:op} that $\m{N}_{12}^{*} \m{N}_{12} = \m{C}/\m{A} = \m{B}/\m{A}$.
Therefore, it suffices to prove that there is some $\m{N}_{12}: \c{H}_{2} \rightarrow \c{H}_{1}$ such that $\c{R}(\m{N}_{12}) \subset \c{N} (\m{B}_{11}^{1/2} \m{A}_{11}^{1/2})$ and $\m{N}_{12}^{*} \m{N}_{12} = \m{B}/\m{A}$. Given that
\begin{align*}
    \rk (\m{B} / \m{A}) = \rk [(\m{B} / \m{A})^{1/2}] &= \dim [\c{R}(\m{B}^{1/2}) \cap \c{H}_{2}] \\
    &\le \dim [\c{N}(\m{B}^{1/2} \m{A}^{1/2}) \cap \c{H}_{1}] = \dim \c{N} (\m{B}_{11}^{1/2} \m{A}_{11}^{1/2}),
\end{align*}
due to \cref{lem:Schur:op}, there is a partial isometry $\m{U}_{12}$ with initial space $\overline{\c{R}(\m{B}/\m{A})}$, and final space included in $\c{N}(\m{B}_{11}^{1/2} \m{A}_{11}^{1/2})$. Then, $\m{N}_{12} = \m{U}_{12} (\m{B}/\m{A})^{1/2}$ is a desired solution, and any such solution takes this form.
\end{itemize}
\end{proof}

\end{document}
